\subsection{The actual relative deformation equations}
\label{sec:support-relative-presentation}

The number of global equations is fixed by an exact Galois cohomology
calculation.  This calculation concerns the original finite inertia
quotients and the original local ring; it does not require a bound on
the degree of the auxiliary automorphic field.

Retain an original good datum over $\mathbf Q$, its complete
coefficient DVR $O$ with residue field $k$ of odd characteristic
$\ell$, and its base representation
\[
 r_0(g)=\bigl(\operatorname{diag}(\Psi(g),1,\Psi(g)^{-1}),
                  \eta(g)\bigr),
 \qquad \Psi(c)=1,\quad\eta(c)=-1.
                                                               \tag{RP1}
\]
Here $c$ is complex conjugation, $\Psi=\varepsilon^e\chi$, and
$\chi$ has order prime to $\ell$.  Write $\psi=\overline\Psi$
and $\omega=\overline\varepsilon$.  Let $R_\ell$ be the original
chosen quotient of the framed local deformation ring, carrying its
universal local representation.  It is kept intact, including all
its integral equations.  Its formal torus action is conjugation by
$\operatorname{diag}(t,1,t^{-1})$.
Fix any supplied weighted local presentation
$O[[Y_u]]/J_\ell\simeq R_\ell$ whose coordinates are the prescribed
local entry differences and whose matrix entries have their matrix
weights.  This is the original local presentation condition, with
no assertion about a twelve-variable characteristic-zero chart.

Let $R$ represent the global deformations with fixed multiplier $\eta$,
the prescribed local representation obtained from $R_\ell$, and the
original condition away from $\ell$: inertia acts trivially on
$\ker\eta\cap\ker\Psi$.  The normalizations at $c$ and the second
mark are imposed later.  All these are framed deformations; no
quotient by conjugacy is taken in their definition.

\begin{theorem}[The original relative presentation]
\label{thm:support-relative-presentation}
There are nonnegative integers $a_j$, zero outside
$j\in\{-2,-1,0,1,2\}$, and a homogeneous presentation
\[
 R\simeq
 R_\ell[[X_{j,s}:1\le s\le a_j]]/(F_{j,t}),
 \qquad
 \#\{t:F_{j,t}\text{ has weight }j\}
       =a_j+\begin{cases}1,&|j|\le1,\\0,&|j|=2.\end{cases}
                                                               \tag{RP2}
\]
Each variable has weight $j$ and maps to a difference between an
actual global matrix entry and its base value.  The old local
coordinates and every matrix entry retain their original weights.
Relative to every supplied local presentation above, the presentation
is $O[[Y_u,X_{j,s}]]/(J_\ell,F_{j,t})$, with the same old
coordinates and equations.
Zero equations may be used to reach the displayed counts.  The
result is the original relative-presentation assertion for every
choice of the original local ring \cite{support-relative-presentation-source}.
\end{theorem}

\begin{proof}
Choose a finite set $\Sigma$ of rational places containing $\ell$,
the ramification of $\eta$, and the conductor of $\chi$.
The global representations factor through $G_{\mathbf Q,\Sigma}$.
For $q\in\Sigma\setminus\{\ell\}$ put
\[
             N_q=I_q\cap\ker\overline r_0.
\]
On inertia away from $\ell$, $\Psi=\chi$ has prime-to-$\ell$
order.  Reduction is injective on both its image and that of $\eta$.
Thus $N_q=I_q\cap\ker r_0$, and these are exactly the inertia
subgroups killed in the original deformation problem.  The group
$I_q/N_q$ has order prime to $\ell$, and $D_q/I_q$ is procyclic.

The deformation condition of factoring through $D_q/N_q$ is formally
smooth.  To see this over a small extension, first lift the finite
inertia representation: averaging makes its positive cohomology
zero, so it lifts uniquely up to conjugation congruent to one.
Its lift and its Frobenius conjugate can therefore be identified by
such a conjugation.  A lift of Frobenius then completes the
representation.  Equivalently, $D_q/N_q$ has $\ell$-cohomological
dimension at most one.  This argument preserves the literal inertia
quotient; it does not replace it by all of $I_q$.

Let $M_j$ be the weight-$j$ part of the residual polarized adjoint.
It has dimensions $1,2,3,2,1$ at weights $-2,-1,0,1,2$.
At $c$ its action is $X\mapsto-JX^{\mathsf t}J$.  Consequently
\[
 b_j:=\dim_kM_j^{c=1}
       =\begin{cases}1,&j=-1,0,1,\\0,&j=-2,2.\end{cases}
                                                               \tag{RP3}
\]
For weight zero the fixed line is
$\operatorname{diag}(1,0,-1)$; for each of weights $1,-1$ it is
the negative combination of the two reflected entries.  The
outer-corner entries at weights $2,-2$ have eigenvalue $-1$.

The original kernel-separation condition provides an element $h$
such that $\eta(h)=-1$ and $\psi(h)=\omega(h)=1$.  Since
$\psi(c)=1$ and $\omega(c)=-1$, the element $ch$ acts trivially
on every $M_j$ while $\omega(ch)=-1$.  It follows that
\[
                  H^0(\mathbf Q,M_j^\vee(1))=0.               \tag{RP4}
\]
This is an exact consequence of the original good-data condition;
it is not a residual largeness hypothesis.

We spell out the relative complex, so that the framing correction is
visible.  Use continuous cochains, and write
\[
 \begin{aligned}
 Q_q^\bullet&=
 C^\bullet(D_q,M_j)/C^\bullet(D_q/N_q,M_j),
 \\
 C_{\rm rel}^\bullet&=
 \operatorname{Cone}\left(
 C^\bullet(G_{\mathbf Q,\Sigma},M_j)
 \longrightarrow C^\bullet(G_{\mathbf Q_\ell},M_j)
                  \oplus\bigoplus_{q\ne\ell}Q_q^\bullet
                         \right)[-1].
 \end{aligned}                                                \tag{RP5}
\]
The inflation map in the first expression is injective on cochains.
The original maps from local absolute Galois groups are used; no
embedding of a local group into $G_{\mathbf Q,\Sigma}$ is assumed.
Degree zero at $\ell$ is retained.  A relative first cohomology class
can be represented by a global cocycle and a local degree-zero
element.  Subtracting that element as a global coboundary makes the
local degree-zero component zero.  It then consists exactly of the
global framed cocycles which vanish at $\ell$ and on every $N_q$.
Thus
\[
 a_j=\dim_kH^1(C_{\rm rel}^\bullet)
\]
is the number of minimal new relative matrix coordinates of weight $j$.

All these cohomology groups are finite.  Finite Poitou--Tate duality
and (RP4) make \cite{AuxTowerDuality}
\[
 H^2(G_{\mathbf Q,\Sigma},M_j)\longrightarrow
 H^2(G_{\mathbf Q_\ell},M_j)
              \oplus\bigoplus_{q\ne\ell}H^2(D_q,M_j)
\]
surjective.  Hence $H^3(C_{\rm rel}^\bullet)=0$.
Also $H^0(C_{\rm rel}^\bullet)=0$, because restriction on invariants
to the distinguished local group is injective.

The local Euler characteristic is $-\dim_kM_j$ at $\ell$ and
zero at the other finite places.  Since $I_q/N_q$ has order prime
to $\ell$, inflation and the cohomology of a procyclic group give
\[
 H^i(D_q/N_q,M_j)=0\ (i\ge2),\qquad
 \dim H^1(D_q/N_q,M_j)=\dim H^0(D_q,M_j).
\]
Thus $Q_q^\bullet$ has Euler characteristic zero.  The global
Euler characteristic over $\mathbf Q$ is $-\dim_kM_j^{c=-1}$.
Taking Euler characteristics in (RP5) now gives the exact identity
\[
 \dim_kH^2(C_{\rm rel}^\bullet)-a_j
       =\dim_kM_j-\dim_kM_j^{c=-1}=b_j.                        \tag{RP6}
\]

For clarity we give the presentation argument, including its integral
relation bound.  Local and global Galois groups with the stated finite
ramification have finite first cohomology for every finite residual
module.  Here is an explicit consequence for representability.
Let $H$ be the kernel of the residual representation.  A lift sends
$H$ into the pro-$\ell$ group of matrices congruent to one, so factors
on $H$ through its maximal pro-$\ell$ quotient $P$.
The group $P$ is finitely generated: its Frattini quotient is dual to
the finite space $H^1(H,\mathbf F_\ell)$.  The kernel of $H\to P$
is characteristic in $H$ and hence normal in the whole group.
The resulting quotient is an extension of the finitely generated
group $P$ by the finite residual image, so is finitely generated.

Choose finitely many topological generators of that quotient.  Their
matrix entries give finitely many formal coordinates with the prescribed
residual values and multiplier.  Invertibility is automatic since the
residual determinants are units.  In their power-series ring impose
the entries of every group relation, using the closed ideal they
generate.  Profinite words are evaluated by continuity in each finite
coefficient quotient.  This constructs the universal framed ring as
a complete Noetherian quotient of a power-series ring over $O$.
It also proves uniqueness from the generating matrices.  The same
argument applies locally.  Imposing the original local quotient and
the specified inertia kernels adds their literal closed equations.

Matrix-entry functionals separate framed relative tangent vectors.
Choose a basis of their dual space, weight by weight, from those
functionals.  The corresponding actual entry differences, together
with the local coordinates, topologically generate $R$.  This gives
a surjection
\[
                   A=R_\ell[[X_{j,s}]]\longrightarrow R
\]
with $a_j$ new variables of weight $j$.  Its kernel $J$ is
homogeneous for the formal torus action.  Weight components of a
formal series are preserved exactly, over $O$ and over $k$; in
particular no use is made of the possibly small group $k^\times$.
Choose homogeneous lifts of a basis of $J/\mathfrak m_AJ$.
They generate the entire closed ideal $J$ by Nakayama and completeness.

Their number in weight $j$ is at most
$\dim_kH^2(C_{\rm rel}^\bullet)$.  Here is the usual obstruction
injection in this specific setting.  A linear functional on
$J/\mathfrak m_AJ$ pushes out
$A/\mathfrak m_AJ\to R$ to a small extension with one-dimensional
kernel.  Continuous matrix lifts exist also for this complete small
extension.  Its finite kernel has zero intersection with some open
additive subgroup of the source.  That subgroup maps homeomorphically
onto an open subgroup of the target; choosing finitely many coset
representatives extends its inverse to a continuous set-theoretic
section.  Invertible residual matrices then have continuous invertible
lifts.  Their product
defects form a global two-cocycle.  Comparing at $\ell$ with the
given $R_\ell$ representation supplies its local one-cochain;
at the other places compare with the formally smooth lifts through
$D_q/N_q$.  Modulo the allowed local cochains these are exactly a
degree-two cocycle in (RP5).  Changing the chosen lifts changes it by
a coboundary.  The obstruction is zero precisely when a compatible
global representation lifts.  The construction is linear in the
functional.  For a relation functional supported in weight $j$, give
the one-dimensional kernel a generator of weight $j$.  Conjugation
multiplies a matrix entry in position $(i,i')$ by
$t^{w_i-w_{i'}}$.  Its product defect has coefficients in that kernel,
so only positions with $w_i-w_{i'}=j$ can occur.  The obstruction
therefore takes values in $M_j$, with the same $j$ as the relation;
no opposite weight is introduced by dualizing the relation space.

If this obstruction is zero, representability gives a section of
the small extension over $R_\ell$.  Compose the section with
$A\to R$.  Every $X_{j,s}$ changes by an element of the small kernel.
Minimality of the relative variables says
\[
                 J\subseteq\mathfrak m_{R_\ell}A+(X)^2.
\]
Taylor expansion therefore shows that this change acts trivially on
the image of $J/\mathfrak m_AJ$: its linear coefficients lie in
$\mathfrak m_A$, which annihilates the small kernel.  The composite
also kills $J$, so the original functional was zero.  This proves
the injection from the dual relation space into the obstruction
space without a flatness assumption or saturation of $J$.

By (RP6) at most $a_j+b_j$ equations of weight $j$ are required.
Adding zero equations reaches that number exactly.  For weights
outside the displayed five, the tangent and obstruction modules
are zero.  Finally lift each homogeneous equation from
$R_\ell[[X]]$ to $O[[Y,X]]$ and take its same-weight component;
the homogeneous local presentation makes this a lift of the original
equation.  The kernel of $O[[Y,X]]\to R_\ell[[X]]$ is exactly
the extended ideal $J_\ell$, so these lifts together with the unchanged
old equations generate precisely the global kernel.  The inclusion
of old coordinates is the specified local map, and every new
coordinate is still its chosen actual entry difference.
All coordinates and local data have therefore been preserved,
which proves (RP2).
\end{proof}

The theorem supplies the global equations used in the original graded
model.  Its separate local input is the twelve-dimensional ordinary
chart at the base representation, with its prescribed five weight
multiplicities.  The integral lower comparison must still use that
same chart and the original strict local conditions.  Neither assertion
is inferred merely from the equation count.
